% Cote 119: offprints, reports, courses, CV and letters (typed and printed part): archive cote 119, site page 218 (PDF page 219).
% typescript. Draft compared: none.
% AI-transcribed by Claude Opus 5.5 (claude-opus-5-5, adaptive thinking, effort high), Message Batches API, batch r119_full (msgbatch_01DwZNm5Y7YsKBWyAN4sFs8w), 2026-11-13; not yet checked by eye.
\scanpage{119/218}
\mspage{. 5 .}

de mes résultats en géométrie algébrique se spécialisent
en des résultats en algèbre pure, comme la relation
$K(A[t]) \simeq K(A)$ , où $A$ est un anneau. Mises à part
ces retombées, on peut signaler les contributions ci-dessous.

a) \uline{Algèbre catégorique} : En fait, de façon continuelle
depuis 1953, je me suis senti dans l'obligation, au
fur et à mesure des besoins, de développer une panoplie
catégorique toujours insuffisante. La plupart des
résultats et des notions ainsi introduites se trouvent
développés un peu partout dans [15,16], notamment dans
le premier exposé de SGA 4. Il ne peut être question
de passer en revue ici même sommairement les notions
qui sont ainsi entrées dans l'usage courant. Signalons
seulement ici le langage des \uline{univers} (pour éliminer
des difficultés logiques dans la manipulation intensive
des catégories), et celui de la \uline{descente} (développé de
façon systématique par Giraud [$*$]) .

\uline{Références} : J. Giraud, P. Gabriel.

b) \uline{Algèbre commutative} : Dans le langage géométrique
des ``schémas'', l'algèbre commutative peut être considérée
comme étant, essentiellement, l'étude locale des schémas.
C'est ainsi que [15], et notamment le Chap. IV de cet
ouvrage, contient de très nombreux résultats nouveaux
d'algèbre commutative, dont il ne peut être question
ici d'énumérer même les plus couramment utilisés. Notons
seulement ici, en algèbre locale, la notion d'anneau
\uline{excellent} et ses propriétés de permanence (dont l'absence
constituait sans doute la lacune la plus marquante de
l'ouvrage de M. Nagata sur les anneaux locaux).

\uline{Références} : M. Nagata, P. Samuel, M. Raynaud, O. Zariski.

\margin{218}
